\chapter{Methodology}
\label{sec:chapterMethodology}

\section{Methodology}
\label{sec:methodology}

% ──────────────────────────────────────────────────────────────────────────
\subsection{Research Objective and Experimental Design}
\label{sec:objective}
% ──────────────────────────────────────────────────────────────────────────

This thesis investigates fraud detection in financial transactions under extreme class imbalance. The work is structured around four complementary research contributions:

\begin{enumerate}
  \item \textbf{Unified model comparison.} Compare classical Machine Learning (ML) models and a Transformer-based tabular model across distinct fraud domains using a single, leakage-free evaluation protocol.

  \item \textbf{Threshold selection as a first-class experimental variable.} Systematically study how the choice of decision threshold strategy affects model rankings and operational metrics, addressing a gap in the fraud detection literature where threshold selection is often implicit or applied on the test set (introducing leakage).

  \item \textbf{Full factorial analysis of imbalance handling.} Evaluate $s$ balancing strategies $\times$ $m$ model families $\times$ $d$ datasets in a structured factorial design, producing a practitioner-oriented recommendation framework for when and for which model each strategy helps.

  \item \textbf{Cross-domain generalization under distribution shift.} Assess model robustness by training on one dataset and evaluating on related datasets with different distributional characteristics, simulating a realistic deployment scenario where a fraud model trained in one context is applied in another.
\end{enumerate}

In addition, the study analyses performance--interpretability trade-offs through SHAP-based explanations.

% ──────────────────────────────────────────────────────────────────────────
\subsection{Datasets}
\label{sec:datasets}
% ──────────────────────────────────────────────────────────────────────────

\paragraph{Credit Card Fraud Detection 2013 (ULB).}
The ULB dataset contains European credit card transactions with anonymised features (PCA components V1--V28, plus \texttt{Time} and \texttt{Amount}) and a binary target indicating fraud. The dataset exhibits extreme class imbalance ($\approx$0.17\% fraud).

\paragraph{Bank Account Fraud (BAF) Suite (NeurIPS 2022).}
The BAF suite contains a Base dataset and five variants (Variant~I--V) designed to stress-test model robustness under distribution shifts and bias-related conditions. Each dataset comprises $\approx$1\,000\,000 instances with mixed numerical and categorical variables (5 categorical features: \texttt{payment\_type}, \texttt{employment\_status}, \texttt{housing\_status}, \texttt{source}, \texttt{device\_os}) and a binary label \texttt{fraud\_bool} ($\approx$1.1\% fraud). The six datasets share a common core of 32 features; Variants~III and~V include two additional synthetic columns (\texttt{x1}, \texttt{x2}), which are dropped in cross-domain experiments to maintain a consistent feature space across all six datasets.

\paragraph{Cross-dataset compatibility note.}
The ULB and BAF datasets have incompatible feature spaces (PCA-anonymised vs.\ real-world features). Therefore, cross-dataset transfer between ULB and BAF is not feasible. Cross-domain generalization is studied \emph{within} the BAF suite, where the shared feature space isolates the effect of distribution shift from feature engineering differences.

% ──────────────────────────────────────────────────────────────────────────
\subsection{Models}
\label{sec:models}
% ──────────────────────────────────────────────────────────────────────────

\subsubsection{Classical ML Models}
The classical baseline includes:
\begin{itemize}
  \item Logistic Regression (LR)
  \item Random Forest (RF)
  \item LightGBM (LGBM)
  \item CatBoost (CB)
  \item One-Class SVM (OCSVM), evaluated as an anomaly-detection baseline
\end{itemize}

\subsubsection{Transformer-based Tabular Model}
A Feature Tokenizer Transformer (FT-Transformer) is used as the main deep learning model for tabular data, modelling feature interactions via self-attention. Beyond performance comparison, the FT-Transformer is used to test whether attention-based architectures are inherently more or less robust to class imbalance than tree-based models---a question that remains open in the literature.

% ──────────────────────────────────────────────────────────────────────────
\subsection{Imbalance Handling Strategies}
\label{sec:imbalance_strategies}
% ──────────────────────────────────────────────────────────────────────────

To address extreme class imbalance, the following strategies are evaluated for supervised models (LR, RF, LGBM, CB, FT-Transformer):
\begin{itemize}
  \item None (baseline---no resampling, no cost-sensitive adjustment)
  \item Random Under-Sampling (RUS)
  \item Random Over-Sampling (ROS)
  \item SMOTE
  \item SMOTE + Tomek Links
  \item SMOTEENN
  \item Class weighting (cost-sensitive learning)
\end{itemize}

OCSVM is evaluated separately (no resampling or class weighting), given its one-class anomaly detection paradigm.

The full factorial design yields $7 \text{ strategies} \times 5 \text{ supervised models} = 35$ configurations per dataset (plus OCSVM), enabling a structured analysis of the \emph{strategy $\times$ model $\times$ dataset} interaction.

% ──────────────────────────────────────────────────────────────────────────
\subsection{Leakage-Free Evaluation Protocol}
\label{sec:protocol}
% ──────────────────────────────────────────────────────────────────────────

\subsubsection{Outer Split (Holdout Test)}
For each dataset, an outer stratified split creates a holdout test set (80\% train, 20\% test). The holdout test set is never used for preprocessing fitting, resampling, hyperparameter selection, threshold selection, or interpretability calibration.

\subsubsection{Inner Cross-Validation (Model Selection and Thresholding)}
All selection procedures occur exclusively within the training partition via stratified $k$-fold cross-validation ($k=5$). For each fold:
\begin{enumerate}
  \item Fit preprocessing on the fold's training subset only.
  \item Apply the imbalance strategy only on the fold's training subset.
  \item Train the model.
  \item Evaluate on the fold's validation subset: compute PR-AUC and obtain predicted scores.
  \item Select a decision threshold $\tau$ by maximising the chosen criterion (see Section~\ref{sec:threshold_study}) on the validation subset.
  \item Using $\tau$, compute $F_1$, $F_2$, and confusion matrix counts $\{TN, FP, FN, TP\}$.
\end{enumerate}

\subsubsection{Final Training and Test Evaluation}
For each (model, strategy, dataset), the final configuration is trained on the full training partition. The decision threshold $\tau$ applied to the holdout test set is derived from inner validation (median threshold across folds) to avoid test-set leakage.

% ──────────────────────────────────────────────────────────────────────────
\subsection{Threshold Selection Study}
\label{sec:threshold_study}
% ──────────────────────────────────────────────────────────────────────────

Most fraud detection studies either use a fixed threshold of 0.5 or optimise the threshold on the test set, which constitutes data leakage. This thesis treats threshold selection as a first-class experimental variable and systematically compares the following strategies:

\begin{enumerate}
  \item \textbf{Fixed $\tau = 0.5$} (na\"ive baseline used in most classification literature).
  \item \textbf{$\tau$ maximising $F_1$} on the validation set (balanced precision--recall trade-off).
  \item \textbf{$\tau$ maximising $F_2$} on the validation set (recall-weighted, more relevant for fraud detection where missing a fraud is costlier than a false alert).
  \item \textbf{$\tau$ at a desired precision operating point} (e.g., $\text{Precision} \geq 0.5$), simulating a constraint-driven deployment where the investigation team has a maximum tolerable false-alert rate.
\end{enumerate}

For each strategy, $\tau$ is selected \emph{exclusively on the validation folds} (never on the test set), and the median across folds is applied to the holdout test set. The primary analysis uses Strategy~3 ($F_2$-maximising), while the threshold study reports all four to quantify how model rankings, $F_1$, $F_2$, and operational metrics change across strategies.

% ──────────────────────────────────────────────────────────────────────────
\subsection{Hyperparameter Tuning Strategy}
\label{sec:tuning}
% ──────────────────────────────────────────────────────────────────────────

Hyperparameter tuning is performed only for the baseline setting (\texttt{strategy=None}) to ensure feasibility and fairness across the imbalance strategy study. The best hyperparameters found under the baseline are reused for all resampling strategies within the same model and dataset, isolating the effect of the balancing technique from confounding model capacity changes.

For the BAF suite, tuning is performed on the Base dataset and the best hyperparameters are reused across the variants, keeping model capacity constant while testing robustness under dataset shifts.

% ──────────────────────────────────────────────────────────────────────────
\subsection{Cross-Domain Generalization}
\label{sec:cross_domain}
% ──────────────────────────────────────────────────────────────────────────

To assess how well fraud detection models generalize across domains, we study cross-domain transfer within the BAF suite, where all six datasets (Base + Variants~I--V) share a common core feature space (32~columns) but differ in distributional characteristics (covariate shift, label shift, or bias conditions). The extra columns present in Variants~III and~V (\texttt{x1}, \texttt{x2}) are excluded so that the model trained on Base can be applied directly to all Variants without feature mismatch.

The protocol is as follows:
\begin{enumerate}
  \item \textbf{In-domain baseline:} For each BAF dataset, train and evaluate on the same dataset (standard holdout protocol) to establish upper-bound performance.
  \item \textbf{Cross-domain transfer:} Train on BAF~Base (using the full training partition with the best configuration from tuning), then evaluate on the holdout test set of each Variant~I--V \emph{without retraining or re-tuning}.
  \item \textbf{Performance degradation:} Report the absolute and relative drop in PR-AUC, $F_1$, and $F_2$ between in-domain and cross-domain to quantify robustness to distribution shift.
  \item \textbf{Threshold transfer:} The decision threshold $\tau$ obtained from Base validation is applied unchanged to the Variants, testing whether threshold calibration transfers across domains.
\end{enumerate}

This simulates a realistic deployment scenario: a fraud model trained at one institution is deployed at another, or a model trained on historical data encounters distributional drift over time.

\paragraph{Note on ULB $\leftrightarrow$ BAF.}
Cross-dataset transfer between ULB~2013 and the BAF suite is not feasible due to incompatible feature spaces (PCA-anonymised vs.\ real-world features). This limitation is acknowledged; the BAF internal transfer setting provides a controlled environment to study generalization under distribution shift with matched features.

% ──────────────────────────────────────────────────────────────────────────
\subsection{Evaluation Metrics}
\label{sec:metrics}
% ──────────────────────────────────────────────────────────────────────────

\paragraph{PR-AUC.}
The primary metric is the area under the Precision--Recall curve (PR-AUC, computed via \texttt{average\_precision\_score}), which is appropriate under extreme imbalance and is threshold-independent.

\paragraph{$F_1$ and $F_2$.}
Threshold-dependent metrics computed using $\tau$ selected on validation. $F_2$ weights recall higher than precision ($\beta=2$) and is used as the operating-point selection criterion.

\paragraph{Confusion matrix.}
We report $\{TN, FP, FN, TP\}$ to quantify operational trade-offs.

\paragraph{Operational workload.}
We report alert rate ($= \frac{TP+FP}{N}$) and the ratio $\frac{FP}{TP}$ as proxies for human review workload.

\paragraph{Calibration.}
Brier score is reported for models producing probability outputs, assessing calibration quality.

\paragraph{Workload-at-$k$.}
Precision@$k$ and Recall@$k$ (top~0.5\% of scored transactions) simulate an investigator reviewing only the $k$ highest-risk alerts.

\paragraph{Confidence intervals.}
Bootstrap 95\% confidence intervals (1\,000 iterations) are computed on the holdout test set for PR-AUC and $F_2$.

% ──────────────────────────────────────────────────────────────────────────
\subsection{Computational Cost and Reproducibility}
\label{sec:reproducibility}
% ──────────────────────────────────────────────────────────────────────────

Each run logs the full configuration (dataset identifier and SHA-256 hash, split seed, model parameters, preprocessing specification, imbalance strategy, threshold rule), training and inference time, and hardware/software context (Python version, platform). Trained model artefacts and evaluation artefacts (metrics JSON files, PR-curve data, confusion matrix plots) are stored to enable full reproducibility.

% ──────────────────────────────────────────────────────────────────────────
\subsection{Interpretability}
\label{sec:interpretability}
% ──────────────────────────────────────────────────────────────────────────

Interpretability is analysed via SHAP for the top-performing classical models and the FT-Transformer:
\begin{itemize}
  \item \textbf{Global explanations:} Feature importance summaries (mean $|\text{SHAP}|$).
  \item \textbf{Local explanations:} Instance-level explanations for representative TP, FP, and FN cases.
  \item \textbf{Cross-model consistency:} Whether the top-$k$ most important features are consistent across model families and across BAF variants.
\end{itemize}

Due to the PCA-anonymised features in ULB~2013, interpretability analysis is expected to be more actionable on the BAF suite, where features have semantic meaning.
